\chapter{等价断言正例}
\section{结果分析}
方法 A 与方法 B 在当前任务上统计等价。
\section{本章小结}
本章完成比较。

\chapter{等价断言负例}
\section{结果分析}
该目标函数可经等价变换写成约束形式。
\section{本章小结}
本章完成推导。

\chapter{因果归因重大问题}
\section{结果分析}
完整方法的性能提升主要归因于门控模块。
\section{本章小结}
本章完成比较。

\chapter{因果归因局部证据}
\section{结果分析}
完整方法的性能提升主要归因于门控模块。

消融实验中逐项移除门控模块后，RMSE 增大 0.08。
\section{本章小结}
本章完成比较。

\chapter{因果归因章级证据}
\section{结果分析}
消融实验记录了模块 A 的受控对比。

第一段补充数据划分。

第二段补充评价协议。

完整方法的性能提升主要归因于模块 B。
\section{本章小结}
本章完成比较。

\chapter{因果一致性负例}
\section{结果分析}
当前误差变化与门控结构一致，并支持二者存在关联。
\section{本章小结}
本章完成比较。

\chapter{次优比较正例}
\section{结果分析}
表\ref{tab:main}汇总各模型结果。
方法 A 的 RMSE 最低。
方法 B 作为基线参与比较。
工况甲数据用于总体评价。
工况乙数据用于边界评价。
工况丙数据用于稳健性评价。
工况丁数据用于迁移评价。
所有方法采用相同协议。
\section{本章小结}
本章完成比较。

\chapter{次优比较负例}
\section{结果分析}
各模型在当前数据块完成比较。
方法 A 的 RMSE 最低。
方法 B 作为基线参与比较。
工况甲数据用于总体评价。
工况乙数据用于边界评价。
工况丙数据用于稳健性评价。
工况丁数据用于迁移评价。
所有方法采用相同协议。
\section{本章小结}
本章完成比较。

\chapter{浅层图形正例}
\section{结果分析}
图\ref{fig:curve}表明方法 A 的曲线更加贴合。
\section{本章小结}
本章完成比较。

\chapter{浅层图形负例}
\section{结果分析}
图\ref{fig:curve}表明方法 A 的曲线更加贴合，RMSE 降低至 0.12。
\section{本章小结}
本章完成比较。

\chapter{箱线词汇正例}
\section{结果分析}
箱线图显示方法 A 的箱体更紧凑。

该结果说明整体波动较小。
\section{本章小结}
本章完成比较。

\chapter{箱线词汇负例}
\section{结果分析}
箱线图显示方法 A 的误差分布更集中。

其中位数和四分位距均较小，上须未明显增长。
\section{本章小结}
本章完成比较。

\chapter{全称断言正例}
\section{结果分析}
方法 A 在所有指标上均优于基线。
\section{本章小结}
本章完成比较。

\chapter{局部比较负例}
\section{结果分析}
方法 A 仅在该子集取得局部最优。
\section{本章小结}
本章完成比较。

\chapter{阶段混用正例}
\section{结果分析}
选定集的 KS 和 W1 均处于对应界限内。

生成样本的联合分布差异较小。
\section{本章小结}
本章完成比较。

\chapter{阶段规范声明负例}
\section{结果分析}
KS 和 W1 用于评价筛选结果。

选定集不得简称为生成样本并外推到原始候选池。
\section{本章小结}
本章完成比较。

\chapter{实验接口正例}
\section{结果分析}
当前测试块表明方法 A 的 RMSE 较低。

这一观察限定在当前数据与评价协议内。

\chapter{实验接口负例}
\section{结果分析}
当前测试块表明方法 A 的 RMSE 较低。

据此，下一节将评价下游效用。

\chapter{交错流水账正例}
\section{结果分析}
RMSE 为 0.12。原因在于编码模块。MAE 为 0.08。这是因为门控生效。AUC 为 0.91。该差异归因于筛选策略。
\section{本章小结}
本章完成比较。

\chapter{交错流水账负例}
\section{结果分析}
RMSE 为 0.12。MAE 为 0.08。AUC 为 0.91。该结果表明误差较低。原因在于编码模块。这是因为门控生效。
\section{本章小结}
本章完成比较。
